\documentclass[10pt,a4paper]{amsart}
\usepackage{amsmath,amssymb,amsthm,graphicx}
\usepackage{verbatim}
\newenvironment{reference}{\comment}{\endcomment}
%\newenvironment{reference}{}{}
\newenvironment{slogan}{\comment}{\endcomment}
\newenvironment{history}{\comment}{\endcomment}

% For commutative diagrams we use Xy-pic
\usepackage[all]{xy}

% We use 2cell for 2-commutative diagrams.
\xyoption{2cell}
% Theorem environments.
%
\theoremstyle{plain}
\newtheorem{theorem}[subsection]{Theorem}
\newtheorem{proposition}[subsection]{Proposition}
\newtheorem{lemma}[subsection]{Lemma}

\theoremstyle{definition}
\newtheorem{definition}[subsection]{Definition}
\newtheorem{example}[subsection]{Example}
\newtheorem{exercise}[subsection]{Exercise}
\newtheorem{situation}[subsection]{Situation}

\theoremstyle{remark}
\newtheorem{remark}[subsection]{Remark}
\newtheorem{remarks}[subsection]{Remarks}

\numberwithin{equation}{subsection}

% Macros
%
\def\lim{\mathop{\mathrm{lim}}\nolimits}
\def\colim{\mathop{\mathrm{colim}}\nolimits}
\def\Spec{\mathop{\mathrm{Spec}}}
\def\Hom{\mathop{\mathrm{Hom}}\nolimits}
\def\Ext{\mathop{\mathrm{Ext}}\nolimits}
\def\SheafHom{\mathop{\mathcal{H}\!\mathit{om}}\nolimits}
\def\SheafExt{\mathop{\mathcal{E}\!\mathit{xt}}\nolimits}
\def\Sch{\mathit{Sch}}
\def\Mor{\mathop{\mathrm{Mor}}\nolimits}
\def\Ob{\mathop{\mathrm{Ob}}\nolimits}
\def\Sh{\mathop{\mathit{Sh}}\nolimits}
\def\NL{\mathop{N\!L}\nolimits}
\def\CH{\mathop{\mathrm{CH}}\nolimits}
\def\proetale{{pro\text{-}\acute{e}tale}}
\def\etale{{\acute{e}tale}}
\def\QCoh{\mathit{QCoh}}
\def\Ker{\mathop{\mathrm{Ker}}}
\def\Im{\mathop{\mathrm{Im}}}
\def\Coker{\mathop{\mathrm{Coker}}}
\def\Coim{\mathop{\mathrm{Coim}}}

% Boxtimes
%
\DeclareMathSymbol{\boxtimes}{\mathbin}{AMSa}{"02}

%
% Macros for moduli stacks/spaces
%
\def\QCohstack{\mathcal{QC}\!\mathit{oh}}
\def\Cohstack{\mathcal{C}\!\mathit{oh}}
\def\Spacesstack{\mathcal{S}\!\mathit{paces}}
\def\Quotfunctor{\mathrm{Quot}}
\def\Hilbfunctor{\mathrm{Hilb}}
\def\Curvesstack{\mathcal{C}\!\mathit{urves}}
\def\Polarizedstack{\mathcal{P}\!\mathit{olarized}}
\def\Complexesstack{\mathcal{C}\!\mathit{omplexes}}
% \Pic is the operator that assigns to X its picard group, usage \Pic(X)
% \Picardstack_{X/B} denotes the Picard stack of X over B
% \Picardfunctor_{X/B} denotes the Picard functor of X over B
\def\Pic{\mathop{\mathrm{Pic}}\nolimits}
\def\Picardstack{\mathcal{P}\!\mathit{ic}}
\def\Picardfunctor{\mathrm{Pic}}
\def\Deformationcategory{\mathcal{D}\!\mathit{ef}}
\usepackage{cleveref}
\crefname{originalreference}{original reference}{original references}
\Crefname{originalreference}{Original reference}{Original references}
\makeatletter
\newlabel{definition-smooth-morphism}{{06HG}{source}{}{sourcepointer-1}{}}
\newlabel{definition-smooth-morphism@cref}{{[originalreference][0][]06HG}{[1][0][]source}{}{}{}}
\makeatother
\crefname{originalreference}{original reference}{original references}
\Crefname{originalreference}{Original reference}{Original references}
\makeatletter
\newlabel{lemma-factor-small-extension}{{06GE}{source}{}{sourcepointer-2}{}}
\newlabel{lemma-factor-small-extension@cref}{{[originalreference][0][]06GE}{[1][0][]source}{}{}{}}
\makeatother
\title{Smoothness of a morphism of cofibred groupoids can be checked on small extensions}
\author{The Stacks Project Authors}
\date{stacks-project:formal-defos.tex:local-snapshot-20261012; GNU FDL 1.2 or later}
\begin{document}
\maketitle
\noindent\textit{Editorial note.} All delivered mathematical text is the author's own, copied from the pinned Stacks Project chapter file formal-defos.tex; no statement or proof step was rewritten, reordered or removed.\par\medskip
\noindent\textit{Editorial note.} The counted proof is the author's complete argument: the factorization of the surjection into small extensions, the induction step combining the induction hypothesis with the small extension case after pushing an object forward, and the assembly of the morphisms required by the definition of smoothness.\par\medskip
\noindent\textit{Editorial note.} The label lemma-smoothness-small-extensions names the counted statement, and the proof environment immediately following it in the source is delivered verbatim as its proof.\par\medskip
\noindent\textit{Editorial note.} The counted statement is carried by the author's own lemma environment in the chapter source.\par\medskip
\noindent\textit{Editorial note.} Pre-open gate: python operations/agent-authoring/tools/precard\_gate.py --source formal-defos.tex --label lemma-smoothness-small-extensions --start 2132 --end 2175 --local-id 28019 returned RESULT=OK (L28019 free, occupancy FREE, public interval clear). Candidate taken from sources/stacks-formal-defos/candidates-halfA.json.\par\medskip
\noindent\textit{Editorial note.} The statement and proof are copied from the pinned chapter file; the cited inputs are kept as external pointers with their canonical Stacks tags.\par\medskip
\noindent\textit{Allowed standard prerequisites (editorial).}\begin{itemize}
\item Factorization of a surjection into small extensions. Same chapter, the lemma factoring a surjective ring map in the category of Artinian algebras into a finite sequence of small extensions. Supplies the factorization B = B\_n -> ... -> B\_0 = A along which the induction runs.
\item Definition of smoothness for a morphism of cofibred groupoids. Same chapter, the definition of a smooth morphism through the lifting of objects and morphisms along surjections. Supplies the shape of the data to be lifted, namely y over B, x over A and a morphism y -> phi(x).
\end{itemize}\medskip
\ifdefined\hypertarget\hypertarget{sourcepointer-1}{}\fi
\noindent\textit{Source reference.} Original reference 06HG: sources/stacks-formal-defos/pinned-original-source.tex, source lines 2113--2113. Prerequisite pointer invoked at proof line(s) 2136; it names a standard fact of the same chapter that the counted argument uses, and every step of the counted proof built on it is retained verbatim in the delivered body.\par\medskip
\ifdefined\hypertarget\hypertarget{sourcepointer-2}{}\fi
\noindent\textit{Source reference.} Original reference 06GE: sources/stacks-formal-defos/pinned-original-source.tex, source lines 497--497. Prerequisite pointer invoked at proof line(s) 2144; it names a standard fact of the same chapter that the counted argument uses, and every step of the counted proof built on it is retained verbatim in the delivered body.\par\medskip
\begin{lemma}
\label{lemma-smoothness-small-extensions}
Let $\varphi : \mathcal{F} \to \mathcal{G}$ be a morphism of categories
cofibered in groupoids over $\mathcal{C}_\Lambda$.  Then $\varphi$ is smooth
if the condition in Definition \ref{definition-smooth-morphism} is assumed to
hold only for small extensions $B \to A$.
\end{lemma}

\begin{proof}
Let $B \to A$ be a surjective ring map in $\mathcal{C}_\Lambda$.
Let $y \in \Ob(\mathcal{G}(B))$, $x \in \Ob(\mathcal{F}(A))$,
and $y \to \varphi(x)$ be a morphism lying over $B \to A$. By
Lemma \ref{lemma-factor-small-extension} we can factor $B \to A$ into
small extensions $B = B_n \to B_{n-1} \to \ldots \to B_0 = A$.
We argue by induction on $n$. If $n = 1$ the result is true by assumption.
If $n > 1$, then denote $f : B = B_n \to B_{n - 1}$ and denote
$g : B_{n - 1} \to B_0 = A$. Choose a pushforward
$y \to f_* y$ of $y$ along $f$, so that the morphism $y \to \varphi(x)$
factors as $y \to f_* y \to \varphi(x)$. By the induction hypothesis
we can find $x_{n - 1} \to x$ lying over $g : B_{n - 1} \to A$ and
$a : \varphi(x_{n - 1}) \to f_*y$ lying over
$\text{id} : B_{n - 1} \to B_{n - 1}$ such that
$$
\xymatrix{
\varphi(x_{n - 1}) \ar[r]_-a \ar[dr] & f_*y \ar[d] \\
& \varphi(x)
}
$$
commutes. We can apply the assumption to the composition
$y \to \varphi(x_{n - 1})$ of
$y \to f_*y$ with $a^{-1} : f_*y \to \varphi(x_{n - 1})$. We obtain
$x_n \to x_{n - 1}$ lying over $B_n \to B_{n - 1}$ and
$\varphi(x_n) \to y$ lying over  $\text{id} : B_n \to B_n$ so that the diagram
$$
\xymatrix{
\varphi(x_n) \ar[r] \ar[d] & y \ar[d] \\
\varphi(x_{n - 1}) \ar[r]^-a \ar[dr] & f_*y \ar[d] \\
& \varphi(x)
}
$$
commutes. Then the composition $x_n \to x_{n - 1} \to x$ and
$\varphi(x_n) \to y$ are the morphisms required by the definition of
smoothness.
\end{proof}

\end{document}
